\chapter{Développement et Réalisation}

\section{Introduction}
Ce chapitre décrit la réalisation concrète du système de détection
du churn client développé pour \textbf{Tunisie Telecom}. Il
documente ce qui est effectivement implémenté dans le code:
un backend \textbf{Django} principal, une intégration vers une
API \textbf{FastAPI} dédiée au traitement Machine Learning, un
mécanisme \textbf{OTP} pour l'authentification sécurisée, un
moteur de notifications, un cycle de vie complet des
recommandations et un export PDF via
\textbf{pdfkit/ wkhtmltopdf}.

Le développement a suivi une approche pragmatique en six axes,
conformément aux sprints définis au Chapitre 2. Les résultats partiels des tests
fonctionnels sont présentés en fin de chapitre, deux scénarios
restant en cours de finalisation au moment de la rédaction.

\section{Architecture fonctionnelle du projet}
Le projet s'appuie sur une architecture Django modulaire, enrichie
d'une API de prédiction FastAPI et d'une interface utilisateur
responsive. La figure\ref{fig:archi_fonctionnelle} illustre
l'organisation globale des composants.

\begin{figure}[H]
    \centering
\includegraphics[width=0.7\linewidth,height=19cm]{images/architecture fonctionnelle.drawio.png}
    \caption{Architecture fonctionnelle du projet}
    \label{fig:archi_fonctionnelle}
\end{figure}

Les responsabilités sont réparties comme suit:

\begin{itemize}
    \item \textbf{Django} gère l'authentification, le dashboard,
    la gestion des utilisateurs, les notifications et l'export PDF.
    \item \textbf{\texttt{learning.models}} contient le c\oe{}ur
    métier: les entités \texttt{Dataset}, \texttt{ClientChurn}
    et les objets auxiliaires liés aux usages clients.
    \item \textbf{\texttt{dashboard.views}} orchestre l'affichage,
    les KPI, la fiche client, l'historique et la génération de
    rapports.
    \item \textbf{\texttt{core}} contient le moteur de
    notifications, la génération de données mock et l'intégration
    vers l'API de prédiction FastAPI.
    \item \textbf{\texttt{config/settings.py}} centralise la
    configuration, notamment la détection du binaire
    \texttt{wkhtmltopdf} et les paramètres \texttt{pdfkit}.
\end{itemize}

Cette organisation garantit une séparation nette entre la couche
métier et données, la couche applicative et workflows, et la
couche présentation.

\section{Sécurité applicative et gestion des habilitations}

\subsection{Authentification à double facteur (OTP)}

L'accès à la plateforme est protégé par un mécanisme
d'\textbf{authentification à deux facteurs} basé sur un code à
usage unique (\textbf{OTP}). Ce choix répond aux exigences de
confidentialité des données clients de Tunisie Telecom.

Le flux se déroule en trois étapes:

\begin{enumerate}
    \item \textbf{Vérification des identifiants}: l'utilisateur
    saisit son nom d'utilisateur et son mot de passe. Django
    vérifie le statut du compte (\texttt{actif},
    \texttt{en\_attente}, \texttt{bloqué}) ainsi que le compteur
    de tentatives échouées.
    \item \textbf{Envoi du code OTP}: un code à six chiffres est
    généré aléatoirement, stocké en base avec une expiration de
    \textbf{dix minutes}, puis transmis par email via
    \textbf{SendGrid}.
    \item \textbf{Validation du code}: le code saisi est comparé
    à l'entrée en base. S'il est valide, la session Django est
    ouverte et le code est \textbf{immédiatement supprimé}: un
    OTP ne peut être utilisé qu'une seule fois.
\end{enumerate}

Ce mécanisme est implémenté dans \texttt{accounts/views.py},
\texttt{accounts/models.py} et
\texttt{core/otp\_email\_service.py}.

La figure \ref{fig:otpCODE} présente le code view de notre moteur OPT tandis que la figure\ref{fig:otpINTERF} présente l'interface de saisie
du code OTP telle qu'elle s'affiche après la première étape de
connexion .

\begin{figure}[htbp]
    \centering
    \begin{subfigure}{0.47\linewidth}
        \centering
    \includegraphics[width=\linewidth,height=7cm]
        {images/otp_view.png}
        \caption{}
    \label{fig:otpCODE}
    \end{subfigure}
    \hspace{0.03\linewidth}
    \begin{subfigure}{0.47\linewidth}
        \centering
   \includegraphics[width=\linewidth,height=7cm]
        {images/INTERFACE OTP.jpeg}
        \caption{}
    \label{fig:otpINTERF}
    \end{subfigure}
\end{figure}

\subsection{Sécurisation des mots de passe}

Aucun mot de passe n'est stocké en clair dans PostgreSQL. Django
utilise l'algorithme \textbf{PBKDF2} couplé à un sel SHA-256,
qui offre deux garanties:

\begin{itemize}
    \item \textbf{Dérivation itérative}: la fonction de hachage
    est appliquée un grand nombre d'itérations, rendant les
    attaques par force brute extrêmement coûteuses.
    \item \textbf{Salage}: une chaîne aléatoire unique par
    utilisateur garantit que deux mots de passe identiques
    produisent des empreintes distinctes, neutralisant l'usage
    des \textit{rainbow tables}.
\end{itemize}

\subsection{Protection contre les attaques par force brute}

Un mécanisme de verrouillage automatique a été développé: à la
troisième tentative échouée, le statut du compte bascule sur
\texttt{bloqué}. L'accès est alors refusé même en cas de saisie
correcte, jusqu'à l'intervention explicite du Chef d'Agence ou
de l'administrateur compétent.

\subsection{Workflow hiérarchique de validation des comptes}

Notre système distingue cinq rôles: Super Administrateur,
Administrateur de Ville, Chef d'Agence, Agent Marketing et Agent
Commercial. La chaîne de validation respecte strictement le
périmètre de chaque acteur: toute tentative d'accès hors périmètre
retourne une erreur HTTP 403. Des règles d'unicité garantissent
qu'il ne peut exister qu'un seul Administrateur par ville et un
seul Chef d'Agence par agence.

La figure \ref{fig:interface_gestion_compte} présente l'interface de
gestion des comptes de l'une des agences en base.

\begin{figure}[htbp]
    \centering    \includegraphics[width=0.7\linewidth,height=7cm]{images/gestion des compte.jpeg}
    \caption{Interface de gestion des compte}
    \label{fig:interface_gestion_compte}
\end{figure}

\section{ Modélisation des données et pipeline ML}

\subsection{Le modèle \texttt{ClientChurn}}

Le modèle principal est \texttt{learning.models.ClientChurn}. Il
représente un client télécom enrichi de variables statistiques,
comportementales et d'un score de churn. Ses attributs sont
regroupés en plusieurs familles, La figure\ref{fig:code_clientchurn} présente l'extrait de code
de cette classe.

\vspace{4 cm}

\begin{figure}[htbp]
    \centering

    \begin{subfigure}{0.47\linewidth}
        \centering
        \includegraphics[width=\linewidth,height=5cm]
        {images/clientchurncode1.png}
        \caption{}
        \label{fig:clientchurn_a}
    \end{subfigure}
    \hspace{0.03\linewidth}
    \begin{subfigure}{0.47\linewidth}
        \centering
        \includegraphics[width=\linewidth,height=5cm]
        {images/clientchurncode2.png}
        \caption{}
        \label{fig:clientchurn_b}
    \end{subfigure}

    \vspace{0.8cm}

    \begin{subfigure}{0.75\linewidth}
        \centering
        \includegraphics[width=0.7\linewidth,height=5cm]
        {images/clientchurncode3.png}
        \caption{}
        \label{fig:clientchurn_c}
    \end{subfigure}

    \caption{Implémentation progressive du modèle \texttt{ClientChurn}
    : lecture de gauche à droite puis de haut en bas}
    
    \label{fig:code_clientchurn}
\end{figure}

\subsection{Logique de classification du risque}

La classification du risque est basée sur le score de churn calculé
par le modèle ML. Le client est classé à \textbf{risque élevé} si
\texttt{score\_churn}$\geq 0{,}32$, et à \textbf{risque faible}
dans le cas contraire. Cette règle est formalisée ainsi:

\begin{equation}
\text{churn\_predit} =
\begin{cases}
\text{Vrai} & \text{si } \texttt{score\_churn} \geq 0{,}32 \\
\text{Faux} & \text{si } \texttt{score\_churn} < 0{,}32
\end{cases}
\end{equation}

Cette logique est appliquée directement lors de la prédiction par
le service FastAPI ou le pipeline local, et le résultat est stocké
dans le champ \texttt{churn\_predit} de l'entité \texttt{ClientChurn}.

\subsection{Entités associées et explicabilité}

Autour de \texttt{ClientChurn}, plusieurs entités enrichissent
l'analyse:

\begin{itemize}
    \item \texttt{EvenementCDR}: traces d'appel, SMS et données
    mobiles, permettant de calculer la récence et la fréquence
    d'utilisation.
    \item \texttt{InteractionDigital}: usages application web,
    recharges et modifications de profil.
    \item \texttt{DonneeGeospatiale}: couverture réseau, qualité
    locale, position et code postal.
    \item \texttt{Reclamation}: tickets clients avec type, statut
    et délai de résolution.
    \item \texttt{CampagneMarketing} et
    \texttt{InteractionCampagne}: actions marketing et réponses
    client.
    \item \texttt{ShapValeur}: valeur de Shapley de chaque feature
    pour le score de churn du client. La persistance de ces valeurs
    permet d'alimenter la fiche client avec une explication locale
    du score.
\end{itemize}

\subsection{Mode Mock: découplage du développement}

Afin de permettre la progression parallèle du développement
front-end et du pipeline ML, un \textbf{module de génération de
données simulées} a été intégré dans \texttt{core}. Il répond
à trois objectifs:

\begin{itemize}
    \item \textbf{Validation UX/UI}: vérifier que les composants
    visuels réagissent correctement avant l'intégration du modèle
    réel.
    \item \textbf{Tests de robustesse}: simuler des cas extrêmes
    (score $= 0$, score $= 1$, données manquantes) pour vérifier
    la stabilité de la logique métier.
    \item \textbf{Découplage technique}: progresser sur
    l'interface sans attendre la finalisation de la phase
    d'entraînement.
\end{itemize}

\textbf{Restriction de production:} le mode Mock sera strictement désactivé en production (\texttt{MOCK\_MODE=False}).
Un contrôle au démarrage Django devrait lever une exception si
\texttt{MOCK\_MODE=True} est détecté en combinaison avec
\texttt{DEBUG=False}, afin d'empêcher toute activation
accidentelle.

\section{ Interface utilisateur et fonctionnalités métier}

\subsection{Flux d'analyse et intégration FastAPI}

Le workflow d'analyse se déroule dans
\texttt{dashboard.views.accueil} et les pages associées. Django
interroge le service \textbf{FastAPI} pour les prédictions et
les métadonnées du modèle. Les métriques sont chargées depuis
l'artefact
\texttt{churn\_api/fastapi\_artifacts/churn\_metadata\_v1.json}
et exposent: le nom du modèle, l'AUC, la précision, le rappel
et le F1-score.

Django affiche ensuite:

\begin{itemize}
    \item les agrégats de churn (taux global, score moyen);
    \item la répartition des clients par niveau de risque
    (élevé: score $\geq 0{,}32$ / faible: score $< 0{,}32$);
    \item la dernière session d'analyse et l'historique paginé;
    \item les liens de navigation vers la fiche client,
    les notifications et les recommandations.
\end{itemize}

Pour garantir la fiabilité du lancement d'analyse, un
\textbf{mécanisme de verrou de session} empêche les soumissions
multiples: la clé \texttt{analyse\_en\_cours} est positionnée
dans la session Django dès le lancement et libérée dans un bloc
\texttt{finally}, garantissant qu'aucune session ne reste bloquée
en cas d'erreur.

\textit{Note: L'accès au tableau de bord principal est réservé aux chefs d'agence et aux administrateurs. Les agents (commercial et marketing) sont redirigés vers une page d'accueil restreinte.}

La figure\ref{fig:interface_accueil} présente la page d'accueil après une analyse complète.


\begin{figure}[H]
    \centering
\includegraphics[width=0.7\linewidth,height=10cm]{images/accueil.jpeg}
    \caption{Dashboard de monitoring du portefeuille
    clients après analyse}
    \label{fig:interface_accueil}
\end{figure}
\vspace{0.3cm}

\subsection{Cycle de vie du client dans le dashboard}

Un client traverse les étapes suivantes:

\begin{enumerate}
    \item Chargement ou création d'un \texttt{Dataset}
    (import CSV ou génération mock).
    \item Génération ou import des enregistrements
    \texttt{ClientChurn} associés.
    \item Calcul des scores de churn via FastAPI et stockage
    en base.
    \item Consultation globale dans le dashboard: KPIs,
    graphiques, historique.
    \item Ouverture de la fiche client pour analyse détaillée
    (score, SHAP, recommandations).
    \item Export PDF ou action sur recommandation/ notification.
\end{enumerate}

\subsection{Fiche client et explicabilité SHAP}

La fiche client est l'interface de \textbf{diagnostic rapide}.
Elle regroupe le profil contractuel et comportemental du client,
son score de churn et le graphique \textbf{SHAP Waterfall} qui
décompose la contribution de chaque feature à la prédiction:

\begin{itemize}
    \item \textbf{Impact positif (rouge)}: variables augmentant
    la probabilité de churn (récence CDR élevée, retards de
    paiement).
    \item \textbf{Impact négatif (vert)}: facteurs favorisant
    la rétention (ancienneté élevée, nombreux services actifs).
\end{itemize}

Les valeurs SHAP sont persistées dans l'entité
\texttt{ShapValeur} et rechargées à chaque ouverture de la fiche,
sans recalcul. La figure\ref{fig:interface_fiche_shap} présente
cette interface.

\vspace{0.3cm}
\begin{figure}[H]
    \centering
    \includegraphics[width=0.7\linewidth,height=7cm]{images/fiche_client.jpeg}
    \caption{Fiche client avec score de churn et explication SHAP
    Waterfall}
    \label{fig:interface_fiche_shap}
\end{figure}

\subsection{Moteur de recommandations et notifications}

Les recommandations sont créées manuellement par les agents
(commercial ou marketing) depuis la fiche client, ou automatiquement
par le système selon des règles métier. Chaque recommandation suit
un workflow de validation hiérarchique.

\subsubsection{Origines et statuts des recommandations}

Le cycle de vie d'une recommandation dépend de son origine:

\begin{itemize}
    \item \textbf{Recommandation manuelle} (créée par un agent
    depuis la fiche client): créée au statut
    \texttt{pending}, nécessitant une validation par le Chef d'Agence.
\end{itemize}

L'ensemble des transitions possibles est gouverné par
\textbf{plusieurs statuts}:

\begin{enumerate}
    \item \texttt{pending}: statut initial
    exclusif des recommandations manuelles. Le Chef d'Agence
    doit se prononcer avant toute transmission à l'agent.
    \item \texttt{active}: recommandation opérationnelle,
    validée par le chef depuis \texttt{pending}. L'agent reçoit
    une notification in-app.
    \item \texttt{completee}: validation finale par le chef
    après vérification. La recommandation est archivée avec
    traçabilité complète (date, agent, note de résultat).
    \item \texttt{rejected}: recommandation rejetée par le chef
    à n'importe quelle étape, avec motif conservé dans
    \texttt{RejetRecommandation}.
\end{enumerate}

\subsubsection{Flux de rejet via \texttt{RejetRecommandation}}

Lorsqu'un agent considère qu'une recommandation n'est pas
pertinente pour un client donné, il peut soumettre une
\textbf{demande de rejet} motivée: cette demande est
obligatoirement accompagnée d'une explication textuelle et
crée un enregistrement \texttt{RejetRecommandation} en base.
Le Chef d'Agence dispose alors de deux options:

\begin{itemize}
    \item \textbf{Accepter le rejet}: la recommandation est
    \textbf{supprimée définitivement} et la demande passe au
    statut \texttt{accepte}.
    \item \textbf{Refuser le rejet}: l'agent est notifié de
    procéder à l'exécution et la demande passe au statut
    \texttt{refuse}.
\end{itemize}

Ce mécanisme garantit qu'aucune recommandation pertinente ne
peut être ignorée sans contrôle hiérarchique explicite. La
traçabilité des décisions de rejet est conservée de façon
indépendante du cycle de vie de la recommandation elle-même.

La figure\ref{fig:interface_recommandations} présente l'interface
de gestion des recommandations depuis le tableau de bord du
Chef d'Agence.

\begin{figure}[H]
    \centering
    \includegraphics[width=0.7\linewidth,height=13 cm]{images/fiche_onglet2.jpeg}
    \caption{Interface de gestion du cycle de vie des recommandations
    (vue Chef d'Agence)}
    \label{fig:interface_recommandations}
\end{figure}
\vspace{0.3cm}

Le dashboard propose également des pages dédiées aux
notifications: \texttt{notifications.html},
\texttt{notification\_detail.html} et une vue d'archives.
La figure\ref{fig:interface_notifs} présente le centre de
notifications.

\begin{figure}[H]
    \centering
    \includegraphics[width=0.7\linewidth,height=9cm]{images/notif_alerte.jpeg}
    \caption{Centre de notifications in-app}
    \label{fig:interface_notifs}
\end{figure}
\vspace{0.3cm}

\textit{Note: Les notifications par email sont réservées aux événements liés au compte et à l'authentification (création de compte, OTP, réinitialisation de mot de passe, journalisation des connexions). Les notifications liées aux analyses, recommandations et alertes churn ne déclenchent PAS d'envoi d'email, seule la notification in-app est utilisée.}

\subsection{Interface utilisateur responsive}

Le front-end est construit avec \textbf{Bootstrap 5} et un CSS
personnalisé (\texttt{static/css/responsive.css}). Les pages
implémentent plusieurs règles adaptatives:

\begin{itemize}
    \item grilles CSS avec \texttt{grid-template-columns} et
    règles \texttt{@media (max-width: 768px)};
    \item sections basculant en colonne unique sur mobile;
    \item tableaux et cartes adaptables;
    \item boutons et actions accessibles sur petit écran.
\end{itemize}

Les templates principaux sont:
\texttt{dashboard/templates/dashboard/accueil.html},
\texttt{dashboard\_global.html} et \texttt{fiche\_client.html}.

\section{Export et reporting PDF}

L'application propose un seul type d'export PDF à ce jour:

\begin{itemize}
    \item \textbf{Fiche client PDF}: via
    \texttt{dashboard.views.fiche\_client\_pdf}, elle contient
    le profil du client, son score de churn et les recommandations actives.
\end{itemize}

La génération repose sur \textbf{pdfkit/ wkhtmltopdf}, configuré
automatiquement via \texttt{config/settings.py} qui détecte le
binaire selon l'OS. Le flux suit trois étapes:

\begin{enumerate}
    \item \textbf{Rendu du template HTML}: Django collecte les
    données et les injecte dans un template dédié
    (\texttt{fiche\_client\_pdf.html}), distinct du template de
    navigation.
    \item \textbf{Conversion par wkhtmltopdf}: \texttt{pdfkit}
    invoque wkhtmltopdf avec les options A4, marges de 15mm et
    encodage UTF-8. Le rendu est produit entièrement en mémoire
    sans écriture disque intermédiaire.
    \item \textbf{Téléchargement direct}: la réponse HTTP est
    retournée avec le type MIME \texttt{application/pdf} et
    l'en-tête \texttt{Content-Disposition: attachment}.
\end{enumerate}

La figure\ref{fig:interface_rapport_pdf} présente un exemple de
rapport exporté.

\begin{figure}[H]
    \centering
    \includegraphics[width=0.7\linewidth,height=9cm]{images/export.jpeg}
    \caption{Rapport PDF exporté d'une fiche client avec score et recommandations}
    \label{fig:interface_rapport_pdf}
\end{figure}

\section{Déploiement et facilitation de l'installation}


Un soin particulier a été apporté à la \textbf{facilité de
déploiement} sur les postes Windows des équipes de Tunisie
Telecom. L'organisation distingue deux répertoires:
\texttt{bin/} pour les exécutables tiers et \texttt{scripts/}
pour les scripts d'orchestration.

\subsection{Exécutables tiers dans \texttt{bin/}}

Afin d'éliminer toute dépendance à une installation système
préalable, les exécutables tiers sont embarqués directement:

\begin{itemize}
    \item \texttt{bin/wkhtmltopdf/windows/}: binaire wkhtmltopdf
    pour Windows, requis par \texttt{pdfkit} pour la génération
    PDF. Des variantes pour \texttt{linux/} et \texttt{mac/}
    sont incluses pour les environnements de développement
    hétérogènes.
\end{itemize}

\textbf{Note sur les bonnes pratiques:} l'inclusion de binaires
tiers dans le dépôt Git représente un compromis entre portabilité
immédiate et bonnes pratiques de gestion des dépendances. Pour
une version de production pérenne, il serait préférable pour nous, de
référencer ces binaires dans un script d'installation téléchargeant
les sources officielles, afin de garantir leur authenticité et
d'alléger le dépôt.

\subsection{Scripts d'orchestration}

Le répertoire \texttt{scripts/} regroupe les scripts de
démarrage qui lancent séquentiellement les services de
l'application:

\begin{itemize}
    \item \texttt{scripts/start\_all.ps1}: démarre l'intégralité
    des services en une seule commande (serveur Django, service
    FastAPI).
    \item \texttt{scripts/start\_django.ps1}: démarre uniquement
    le serveur Django, utile pour le développement isolé.
\end{itemize}

\subsection{Infrastructure de production envisagée (Docker)}

La conteneurisation de l'application via \textbf{Docker Compose}
est prévue pour le déploiement en production. L'architecture
cible orchestre les services suivants:

\begin{itemize}
    \item \textbf{Nginx}: terminaison SSL et distribution des
    requêtes vers Django et FastAPI.
    \item \textbf{Gunicorn}: serveur WSGI multi-worker pour
    Django.
    \item \textbf{Service FastAPI}: conteneur dédié au pipeline
    ML, indépendant du serveur Django.
    \item \textbf{PostgreSQL}: base de données persistante.
\end{itemize}

Cette configuration garantira un environnement reproductible
et isolé pour chaque composant, et permettra une mise à l'échelle
indépendante du service ML.

\section{Validation et contrôle qualité}

\subsection{Tests unitaires}

Une suite de tests automatisés a été mise en place via
\textbf{Pytest} pour valider les composants critiques:

\begin{itemize}
    \item La logique de hachage PBKDF2 des mots de passe.
    \item Le flux OTP complet (génération, expiration à dix
    minutes et invalidation après usage unique).
    \item Les règles de verrouillage de compte après trois
    tentatives échouées.
    \item Le calcul correct des contributions SHAP pour des
    vecteurs d'entrée connus.
    \item La règle de classification: score $\geq 0{,}32$
    → churn prédit vrai.
\end{itemize}

\subsection{Tests fonctionnels}

\begin{table}[H]
    \small
    \centering
    \begin{tabular}{|p{3.5cm}|p{4cm}|p{3cm}|p{2cm}|}
        \hline
        \textbf{Fonctionnalité} & \textbf{Scénario testé} &
        \textbf{Résultat attendu} & \textbf{Statut} \\
        \hline
        Authentification OTP & Code valide dans la fenêtre
        de 10 min & Redirection accueil & \textbf{Passé} \\
        \hline
        Expiration OTP & Code saisi après 10 min &
        Accès refusé & \textbf{Passé} \\
        \hline
        Verrouillage compte & 3 tentatives échouées &
        Compte bloqué & \textbf{Passé} \\
        \hline
        Validation hors périmètre & Chef valide un agent
        d'une autre agence & Erreur 403 & \textbf{Passé} \\
        \hline
        Seuil de risque & Score $\geq 0{,}32$ &
        Churn prédit vrai & \textbf{Passé} \\
        \hline
        Recommandation manuelle & Agent crée une recommandation &
        Statut \texttt{pending}, chef notifié in-app &
        \textbf{Passé} \\
        \hline
        Demande de rejet & Agent soumet demande motivée &
        \texttt{RejetRecommandation} créé, chef notifié in-app &
        \textbf{Passé} \\
        \hline
        Rejet accepté & Chef accepte le rejet &
        Recommandation supprimée définitivement &
        \textbf{Passé} \\
        \hline
        Workflow validation & Chef valide recommandation manuelle &
        Statut \texttt{active}, agent notifié in-app &
        \textbf{Passé} \\
        \hline
        Anti-doublon notif. & Analyse relancée en moins de 5 min &
        Une seule notification & \textbf{Passé} \\
        \hline
        Verrou de session & Double soumission du pipeline &
        Seconde requête ignorée & \textbf{Passé} \\
        
        \hline
        Export PDF fiche & Clic sur \og{}Exporter\fg{} &
        PDF téléchargé avec SHAP & \textbf{Passé} \\
        \hline
        Responsive (375 px) & Affichage mobile &
        Interface Bootstrap adaptée & \textbf{Passé} \\
        \hline
        Restriction analyse & Agent tente de lancer analyse &
        Accès refusé, message d'erreur & \textbf{Passé} \\
        \hline
        Restriction tableau de bord & Agent accède au tableau de bord &
        Redirection vers page restreinte & \textbf{Passé} \\
        \hline
    \end{tabular}
\end{table}
\begin{table}[H]
    \small
    \centering
    \begin{tabular}{|p{3.5cm}|p{4cm}|p{3cm}|p{2cm}|}
        \hline
        Déploiement Docker & Conteneurisation de l'application &
    Fonctionnalité non encore implémentée &
    \textbf{Non réalisé} \\
        \hline
        Versionnement GitHub & Publication et gestion du dépôt distant &
    Projet non encore publié sur GitHub &
    \textbf{Non réalisé} \\
        \hline
    \end{tabular}
    \caption{Tableau de tests fonctionnels}
    \label{tab:tests_fonctionnels}
\end{table}

Les deux éléments marqués \textit{Non réalisé}, à savoir le déploiement
Docker et la publication du projet sur GitHub, n'ont pas encore
été finalisés au moment de la rédaction de ce rapport. 


\section{Limites et perspectives}

\subsection{Limites techniques identifiées}

\paragraph{Complexité computationnelle de SHAP}
Le calcul des valeurs de Shapley présente une complexité
théorique en $O(2^n)$. Bien que le \texttt{TreeExplainer} de
SHAP optimise ce calcul pour XGBoost, un volume important de
requêtes simultanées pourrait saturer les ressources CPU. Pour
un déploiement à l'échelle nationale, les valeurs SHAP devraient
être calculées de façon asynchrone et mises en cache.

\paragraph{Absence de boucle de rétroaction}
Le moteur de recommandations ne dispose pas d'une \textbf{boucle
de rétroaction} (\textit{Feedback Loop}). Il ne mesure pas encore
si une recommandation a effectivement abouti à la rétention du
client, ce qui empêche le système d'améliorer ses suggestions au
fil du temps.

\paragraph{Mode Mock en production}
Le mode Mock doit être strictement désactivé en production
(\texttt{MOCK\_MODE=False}). Un contrôle au démarrage devrait
lever une exception si ce mode est activé avec
\texttt{DEBUG=False}.

\paragraph{Absence de signature électronique des rapports}
Les rapports PDF sont générés sans signature électronique ni
enregistrement côté serveur. Pour un usage juridiquement
opposable, une évolution vers des rapports horodatés et signés
numériquement serait nécessaire.

\paragraph{Binaires tiers dans le dépôt}
L'inclusion du binaire wkhtmltopdf dans \texttt{bin/} facilite
le déploiement immédiat mais alourdit le dépôt. La migration
vers un script d'installation téléchargeant les sources
officielles est recommandée à terme.

\subsection{Perspectives d'évolution}

\begin{itemize}
    \item \textbf{Conteneurisation Docker}: mettre en \oe{}uvre
    la configuration Docker Compose pour la mise en production.
    \item \textbf{Boucle de rétroaction}: mesurer le taux de
    succès des recommandations pour améliorer le moteur de
    suggestion.
    \item \textbf{Intégration NLP}: analyser les rapports de
    réclamations textuels pour enrichir les features du modèle.
    \item \textbf{Dashboarding longitudinal}: vue analytique
    permettant au Chef d'Agence de suivre l'évolution du taux de
    risque mois par mois.
    \item \textbf{Prédictions en temps réel}: migrer du
    traitement batch vers un pipeline de prédiction en continu
    alimenté par les événements CDR en flux.
\end{itemize}

\section{Conclusion}

Ce chapitre a présenté la réalisation concrète de notre système
de détection du churn client pour Tunisie Telecom. En combinant
la robustesse de \textbf{Django 4.2} pour la logique métier et
la sécurité, le service \textbf{FastAPI} pour le pipeline de
prédiction ML, et l'explicabilité transparente des
\textbf{valeurs SHAP}, nous avons abouti à un outil d'aide à la
décision opérationnel et hiérarchisé.

L'architecture modulaire, les mécanismes de sécurité multicouche
(OTP, PBKDF2, RBAC à cinq niveaux, verrouillage par force brute),
le workflow de validation humaine des recommandations et les
restrictions d'accès par rôle répondent directement aux
exigences fonctionnelles et opérationnelles de Tunisie Telecom.
